\documentclass[10pt]{beamer}

\usetheme{metropolis}
\definecolor{WiLabRed}{RGB}{197,18,48}
\setbeamercolor{frametitle}{fg=white,bg=WiLabRed}
\setbeamercolor{progress bar}{fg=WiLabRed!90}
\setbeamercolor{title separator}{fg=WiLabRed!90}
\setbeamercolor{progress bar in section page}{fg=WiLabRed!90}
\setbeamercolor{background canvas}{bg=white}
\setbeamercolor{alerted text}{fg=WiLabRed!90}

\usepackage{appendixnumberbeamer}

\usepackage{booktabs}
\usepackage[scale=2]{ccicons}

\usepackage{pgfplots}
\usepgfplotslibrary{dateplot}

\usepackage{xspace}
\newcommand{\themename}{\textbf{\textsc{metropolis}}\xspace}

\usepackage{marvosym}
%\usepackage{subfig}
\usepackage{graphicx}\graphicspath{{images/}}
\usepackage{subcaption}
\usepackage[framed]{./support/mcode}
\usepackage{listings}

\title{Lecture 02: SDR Architectures}
\subtitle{\textit{Software Defined Radio for Engineers} (Collins~\textit{et al.}), \textsection{1.5-1.6}}
\date{}
\author{\textbf{Alexander M. Wyglinski, Ph.D.}}
\institute{ \vspace*{1in}\hfill\includegraphics[height=1.125cm]{wilab_logo-A70916.eps} \qquad \includegraphics[height=1.125cm]{WPI_Inst_Prim_FulClr.eps}}
% \titlegraphic{\hfill\includegraphics[height=1.5cm]{logo.pdf}}

% Foot for all slides
\setbeamertemplate{frame footer}{\tiny \copyright~2018 by Alexander Wyglinski. This work is licensed under the Creative Commons Attribution-ShareAlike 4.0 International License. To view a copy of this license, visit http://creativecommons.org/licenses/by-sa/4.0/.}

\begin{document}

%\captionsetup[subfigure]{labelformat=empty}

%%%%%%%%%%%%%%%%%%%%%%%%%%%%%%%%%%%%%%%%%%%%%%%%%%%%%%%%%%

\maketitle




%%%%%%%%%%%%%%%%%%%%%%%%%%%%%%%%%%%%%%%%%%%%%%%%%%%%%%%%%%

\frame
{
  \frametitle{Microelectronics Evolution is Key!}
  \begin{itemize}
    \item Rapid evolution of microelectronics over the past several decades
    \begin{itemize}
        \item {\bf Moore's Law} defines the long-term trend of the number of transistors accommodated on an integrated circuit
        \item Doubling number of transistors per integrated circuit approximately every two years
        \item Affects processing speed and memory
    \end{itemize}
    \item Microelectronics industry has significantly influenced digital communication systems sector
    \begin{itemize}
        \item Digital communication transceivers are becoming more {\it versatile}, {\it powerful}, and {\it portable}
    \end{itemize}
    \item Microelectronics advancements have given rise to {\bf software-defined radio} (SDR) technology
    \begin{itemize}
        \item Baseband radio functions can be {\it entirely} implemented in digital logic and software
    \end{itemize}
  \end{itemize}
}

%%%%%%%%%%%%%%%%%%%%%%%%%%%%%%%%%%%%%%%%%%%%%%%%%%%%%%%%%%

\frame
{
  \frametitle{Types of Microelectronics for SDR Implementation}
  \begin{itemize}
    \item General Purpose Microprocessor
    \begin{itemize}
        \item Very flexible in terms of reconfigurability
        \item Easy to implement new digital communication modules
        \item Not specialized for mathematical computations
        \item Potentially power inefficient
    \end{itemize}
    \item Digital Signal Processor (DSP)
    \begin{itemize}
        \item Specialized for performing mathematical computations
        \item Easy to implement new digital communication modules
        \item Potentially slow for computationally expensive processes
        \item Can be power efficient (e.g., used in cellular telephones)
    \end{itemize}
    \item Field Programmable Gate Array (FPGA)
    \begin{itemize}
        \item Computationally powerful
        \item Power inefficient
        \item Not flexible nor easy to implement new modules
    \end{itemize}
  \end{itemize}
}

%%%%%%%%%%%%%%%%%%%%%%%%%%%%%%%%%%%%%%%%%%%%%%%%%%%%%%%%%%

\frame
{
  \frametitle{SDR Definition}
  \begin{itemize}
    \item A class of reconfigurable/reprogrammable radios whose physical layer characteristics can be significantly modified via software changes
    \begin{itemize}
        \item Capable of implementing different functions at different times on the same platform
        \item Defines in software baseband radio features, e.g., modulation, error correction coding
        \item Possesses some software control over RF front-end operations, e.g., transmission carrier frequency
    \end{itemize}
    \item Baseband radio functionality stored in memory
    \begin{itemize}
        \item Different types of modulation, error correction coding, and other functional blocks are available to the SDR platform
        \item Functional blocks can potentially be changed in real-time
    \end{itemize}
    \item Operating parameters of functional blocks can be adjusted either by human operator or automated process
  \end{itemize}
}

%%%%%%%%%%%%%%%%%%%%%%%%%%%%%%%%%%%%%%%%%%%%%%%%%%%%%%%%%%


\frame
{
  \frametitle{SDR Key Features\footnote{\tiny From {\it Software Radio: A Modern Approach to Radio Engineering} by Jeffrey H. Reed (Prentice Hall, 2002).}}
  \begin{itemize}
    \item {\bf Multifunctionality} -- Possessing the ability to support multiple types of radio functions using the same digital communication system platform
    \item {\bf Global Mobility} -- Transparent operation with different communication networks located in different parts of the world, i.e., not confined to just one standard
    \item {\bf Compactness and Power Efficiency} -- Many communication standards can be supported with just one SDR platform
    \item {\bf Ease of Manufacturing} -- Baseband functions a software problem, not a hardware problem
    \item {\bf Ease of Upgrading} -- Firmware updates can be performed on SDR platform to enable functionality with latest communication standards
  \end{itemize}
}

%%%%%%%%%%%%%%%%%%%%%%%%%%%%%%%%%%%%%%%%%%%%%%%%%%%%%%%%%%


\frame
{
  \frametitle{Digital Communication Transceiver -- Revisted}

  \begin{itemize}
   \item So where do we draw the line between the analog and digital domains in a communication system?
  \end{itemize}

}

%%%%%%%%%%%%%%%%%%%%%%%%%%%%%%%%%%%%%%%%%%%%%%%%%%%%%%%%%%

\frame
{
  \frametitle{History of SDR Development}
  \begin{itemize}
    \item  ``Software-defined radio'' first coined by Joseph Mitola in 1991 although SDR technology was available since the 1970s
    \item First publicly funded SDR development initiative was {\it SpeakEasy I/II} by the U.S. military
    \begin{itemize}
        \item Employed programmable microprocessors for implementing more than ten military communication standards
        \item Transmission carrier frequencies ranged from 2~MHz to 2~GHz
        \item Allowed for upgrades of new functional blocks, such as modulation schemes and coding schemes
        \item SpeakEasy II was the first SDR platform to involve FPGA modules for implementing digital baseband functionality
        \item Physical size of prototype fit in the back of a truck
        \item Read ``SPEAKeasy, the Military Software Radio'' by Upmal and Lackey in {\it IEEE Communications Magazine} (IEEE Press, 1995)
    \end{itemize}
  \end{itemize}
}


%%%%%%%%%%%%%%%%%%%%%%%%%%%%%%%%%%%%%%%%%%%%%%%%%%%%%%%%%%




\end{document}
